%! Logarithmic Function Context and Data Modeling — Unit 2, Lesson 2.14 — Student Packet Cover Sheet
\documentclass[10pt]{article}
\usepackage{precalculus-article}
\usepackage{precalculus-boxes}
\usepackage{ltablex}
\keepXColumns

\begin{document}

% Full-bleed navy banner
\begin{tikzpicture}[remember picture, overlay]
  \fill[navy] (current page.north west)
    rectangle ([yshift=-0.9in]current page.north east);
\end{tikzpicture}
\vspace{-0.8in}

\begin{center}
  {\color{white}\textbf{\LARGE Precalculus}}\\[4pt]
  {\color{skymid}\large\textbf{Unit 2: Exponential and Logarithmic Functions}}\\[6pt]
  {\color{white}\textbf{\large Lesson 2.14 \quad Logarithmic Function Context and Data Modeling}}
\end{center}

\vspace{0.22in}
\namedateperiod
\vspace{0.18in}

\begin{learningtargetbox}
\small
\begin{enumerate}[leftmargin=1.5em, itemsep=5pt]
  \item I can recognize when a data set or context calls for a logarithmic function model.
        \hfill\textit{LO 2.14.A, EK 2.14.A.1}
  \item I can construct a logarithmic model from a proportion and a real zero, or from two
        input-output pairs.
        \hfill\textit{LO 2.14.A, EK 2.14.A.2}
  \item I can apply transformations to $f(x) = a\log_b x$ to fit a logarithmic model to
        a context.
        \hfill\textit{LO 2.14.A, EK 2.14.A.3}
  \item I can use a logarithmic model to predict values for the dependent variable.
        \hfill\textit{LO 2.14.A, EK 2.14.A.6}
\end{enumerate}
\end{learningtargetbox}

\vspace{0.14in}

\begin{tocbox}
\small
\renewcommand{\arraystretch}{1.5}
\begin{tabularx}{\linewidth}{@{} c l X r @{}}
\rowcolor{sky}
\textbf{\#} & \textbf{Component} & \textbf{Description} & \textbf{Score} \\
\midrule
1 & Warm-Up      & Spiral review: table patterns, model identification, real zero & \blank{1.2cm} \\
2 & Guided Notes & When to use a log model; building from two pairs; transformations; predictions & \blank{1.2cm} \\
3 & Group Activity & Species-diversity and decibel logarithmic modeling & \blank{1.2cm} \\
4 & Exit Ticket  & Identify base, write model, predict; evaluate and invert $3\ln(x)$ & \blank{1.2cm} \\
5 & Homework     & Independent practice with logarithmic function modeling & \blank{1.2cm} \\
\midrule
  & \multicolumn{2}{r}{\textbf{Total}} & \blank{1.2cm} \\
\end{tabularx}
\end{tocbox}

\end{document}
